% !TeX root = ../papers/nc-p3.tex
\section{十九次階数二曲線の核曲面と逆線形系}
\label{sec:cofactor-nineteen-surface}
\begin{lemma}\label{lem:cofactor-plane-evaluation}
整な一般被約の階数二曲線の場合、核グラフの解消上の一般入力平面切断を $B$、入力平面八次への射を
$h$、核座標への射を $j$ とする。$j$ が閉埋込みなら
$g(B)=21-d$ である。
\end{lemma}
\begin{proof}
$B$ はBertiniにより滑らかな既約曲線で、$h$ は正規化である。
表示の平面制限から得る純な層 $M_H$ から $h_*j^*\mathcal O(1)$ への評価は一般同型ゆえ単射である。
一般平面は階数一および階数二軌跡の孤立点を避ける。階数三点では繊維の $j$ による
スキーム像は核の一点に入るので、埋込みなら繊維長一である。
階数二の $d$ 点では定理\ref{thm:cofactor-sectional-kernel}の重複度の計算により繊維長二。
射影空間の長さ二スキームでは一次式の評価が全射である。従って全繊維で評価が全射となり、
上の包含は同型である。左辺のEuler数は $4$、右辺は $1+(24-d)-g(B)$ である。
\end{proof}

\begin{theorem}\label{thm:cofactor-nineteen-surface}
$d=19$ なら核曲面は二次曲面で、核写像は双有理である。逆写像の四次元線形系は
滑らかな二次曲面なら母線の交換を除き $(2,3)$、二次錐なら最小解消
$\mathbb F_2$ 上で $2E+5f$（$E^2=-2$, $f$ は母線）である。
どちらでも基点は無限近接点を含めて単純な四点、または重複度二の一点で尽きる。
\end{theorem}
\begin{proof}
定理\ref{thm:cofactor-curve-kernel-only-bezout}と $8n\deg K\le25$ より
$n=1$, $\deg K=2,3$。核曲面の正規化の滑らかなモデル $S_0$ で、固定因子を除いた逆線形系を
$H_0$、核超平面を $L$ とする。基点の重複度を $m_i$ とすれば
$H_0L=5$, $H_0^2=8+\sum m_i^2$。三次ならHodge指数から $H_0^2\le25/3$、
従って基点なしで入力射 $f_0$ が定まる。

滑らかな三次曲面および正規で錐でない三次曲面では、最小解消で $K_{S_0}=-L$。
実際特異な場合は二重点からの射影で有理であり、三重点なら錐になる。
超曲面列から $\chi(\mathcal O_K)=1=\chi(\mathcal O_{S_0})$ で特異点は有理である。
有理Gorenstein曲面特異点の最小解消はcrepantである。従って $H_0(H_0+K_{S_0})=8-5$ が奇数となり矛盾。
三次錐の正規化の母線束は平面三次の正規化上にあり、負切断 $E$ の自乗は $-3$、
$L\equiv E+3f$。$H_0\equiv aE+bf$ として $b=5$, $10a-3a^2=8$ は整数解 $a=2$ だけで、
$H_0E=-1$ は不可能である。

非正規非錐三次曲面の特異曲線は直線 $T$ である。実際特異点二点を結ぶ直線は曲面に含まれ、
非直線曲線の割線の合併は少なくとも平面、非平面曲線では全空間である。
$T=(x_0=x_1=0)$ として式は $za_2(x_0,x_1)+wb_2(x_0,x_1)+c_3(x_0,x_1)=0$。
$a_2,b_2$ は独立（比例なら錐）、三係数に共通零点はない。
$[x_0:x_1]=[u:v]$ の平面の残余直線を集めると、滑らかな直線束 $S_0$ からの射
\[
 (u:v; r,z,w)\longmapsto(ru:rv:z:w),\qquad
 r c_3(u,v)+z a_2(u,v)+w b_2(u,v)=0
\]
を得る。$T$ 上の点の逆像も二次式の根で有限なのでこれは正規化である。
$S_0=\mathbb F_j$ 上で豊富な $L=E+bf$ は $L^2=3$, $LE>0$ より $j=1$, $b=2$。
$H_0=aE+bf$ と置き直すと $a+b=5$, $10a-3a^2=8$ から $H_0=2E+3f$。
これは豊富で $f_0$ は有限、一般切断 $B$ は種数一である。

ここで $B$ の核座標による像は埋込みである。詳細には、$T$ の逆像 $Z$ は $r=0$ で、
直線束の各平面内の式 $za_2+wb_2=0$ が定める。$\gcd(a_2,b_2)=1$ なら
一つの切断が $T$ に二対一、一次共通因子があれば切断と一本の繊維が各々 $T$ に一対一である。
$Z$ の類は $L-f=E+f$（$x_0=ru$ の因子から一繊維を引く）で $H_0 Z=3$。
$T$ の一般点上の二点が常に同じ入力像を持つことはない。既約の二重被覆なら
$Z$ 上の入力射は $T$ を経由し超平面次数が偶数になり、分裂の場合も各枝の次数が等しくなるからである。
よって二入力像が一致する例外は有限で、異なる二入力像を同時に含む平面は各対につき一次元族である。
一般平面は従って $B$ 上で二点を同一視しない。
正規化射の微分も有限個の点以外で単射である。$T$ の外では同型であり、
切断の一般点では $T$ 方向の微分が非零で $r$ 方向は $T$ を出る。
共通根の繊維でも $c_3\ne0$ で、式を $r$ について解くと、底方向の $r$ の一次係数は
一般の $[z:w]$ で非零である。一般 $B$ は有限例外を避け、固有な単射かつ非分岐な射で閉埋込みとなる。
補題\ref{lem:cofactor-plane-evaluation}の種数二に反する。これで三次は全て排除できた。

滑らかな二次曲面では逆系の二次数は和五、積の二倍が八以上なので $(1,4)$ または $(2,3)$。
前者は基点なしで一般切断が有理曲線、核座標が埋込みなので同じ補題に反する。
二次錐では $L=E+2f$, $H_0=aE+5f$ であり $10a-2a^2\ge8$, $5-2a\ge0$ より $a=1,2$。
$a=1$ は基点なし豊富で $f_0$ は有限である。$E$ の入力像は固定された頂点列
$N(x)y_{\rm vertex}=0$ に入る。この列の曲線部分の次数は二以下
（定理\ref{thm:opposite-column-divisor}の四独立二次式の次数評価。曲面因子なら整八次行列式を割る）。
$H_0E=3$ ゆえ $E$ は直線へ三対一で写る。一般平面はその直線と
$C$ の交点を避け、直線の一般点の上の相異なる三点を切り取る。
一般平面八次のこの追加特異点は三枝以上、正規化余長三以上である。
既に $C$ との十九交点が各々余長一以上なので平面八次の種数二十一を超えてしまう。
残る二つの系では $H_0^2=12$。従って $\sum m_i^2=4$ から基点の主張を得る。
\end{proof}

\begin{lemma}\label{lem:symmetric-kernel-graph-reconstruction}
固定座標で上記の核グラフを与える対称二次式行列はスカラー倍を除いて一意である。
パラメータによるグラフ $(x,y)$ があれば、未知対称行列 $P$ に対する線形方程式 $P(x)y=0$ で復元できる。
\end{lemma}
\begin{proof}
実現 $N$ が一つあれば $\operatorname{adj}(N)P$ は対称性により整八次 $D$ 上で消える。
各成分は八次なので $\operatorname{adj}(N)P=(\det N)A$ で $A$ は定数、従って $P=NA$。
$NA=A^tN$ と十成分の独立性から、全対称定数行列 $S$ に対し $SA=A^tS$。
対角行列と各対称行列単位を代入すれば $A$ はスカラーである。
\end{proof}
\path{computations/verify_cofactor_nineteen_surface.py} は各曲面の交点数と残る基点型を整数演算で検算する。
